\scanpage{16}\mspage{15}
\unsure{par|pr} contre la relation $\ll$ n'a pas d'interprétation simple directe, ni en termes des \inserted{\unsure{déployements|déployement}} $\widetilde F\subset\unsure{M|\mathcal{M}}$, ni en termes des multidéployement $\mathrm{Fig}(F)$; ainsi $G\ll F$ n'implique nullement $\mathrm{Fig}_{M}(G)\ll\mathrm{Fig}_{M}(F)$. Aussi l'application ($**$) et la notation $=\mathrm{Fig}_{M}(F)$ sont-\unsure{il|ils} finalement un peu bidon !

En termes de $M$, la structure est donnée par
\[
\begin{array}{|l}\hline
\unsure{\mathcal{F}|\mathfrak{F}|\mathscr{F}}\subset\unsure{\mathcal{P}|\mathscr{P}|\mathfrak{P}|\mathrm{P}}(M)\quad\text{ens de parties de }M\\
\ll\quad\text{relation d'ordre dans }\mathcal{F},\ \text{\uline{impliquée}}\\
\qquad\text{par l'inclusion }(\deleted{\unsure{F}}\text{\unsure{notée|noté} }F\leq G)
\\\hline\end{array}
\tag*{(1.33)}
\]
Ça a l'air moins sympathique du point de vue simplicité des données et d'autant plus qu'il n'y a pas \inserted{à première vue} d'interprétation directe de $\ll$ en termes de \deleted{\unsure{Fig-structure|Fig, structure|notions de structure|notions de théorie}} \inserted{notions ensemblistes} familières. Mais c'est ici que vont servir les (grandes) ombres !

En termes de la \inserted{\unsure{première|premier}} donnée (1.33.\unsure{)|]} \margin{(et avant d'introduire $\ll$)}, le \deleted{\unsure{tout}} premier axiome est \deleted{\illegible{1}} le fait que $\mathcal{F}$ est \deleted{une figure} \inserted{une famille d'}ens. \inserted{\uuline{pré}admissibles} \deleted{\unsure{dans}} $M$ de \margin{attention au ``\uuline{pré}'' !} \deleted{\unsure{Axiome|Axiomes}} support $M$, \deleted{\unsure{d'où}} donc la donnée de \margin{\deleted{\unsure{un ens.|un ensemble|Chaque ens} de \unsure{parties} \unsure{fermées|finies} de $\unsure{F|f|\mathcal{F}|\lbrace}$}} $\mathcal{F}$ \deleted{\unsure{revient à|revient} la donnée d'} \inserted{implique} une relation de préordre dans $M$ (\unsure{ce sera|ce ser.} $\leq$)\unsure{$\big/\!16\!\big/$|$\big/\!6\!\big/$} -- lequel \unsure{sera|ser.}
